% practice_answers.tex -- Module 7 (The Nature of Light): worked answers to the Practice set.
% One practiceans per practiceq in practice.tex, in the same order.
% Needspace values of 22-26 lines keep each short-answer working on one page.
%
% The worked steps (\given ... \checkstep) are used here outside a worked-example box.
% The next line switches off their "SOLUTION" divider line, which is only meant for worked
% examples. It sets an existing flag of hscphys.sty; nothing is redefined.
\global\csname hsc@dividertrue\endcsname
\begin{practiceanswers}

%% ==================================================================
\subsection*{Section I: multiple choice}

\Needspace{13\baselineskip}
\begin{practiceans}
  \correct{C}
  \begin{itemize}
    \item[\xmark] \textbf{A:} a charge at rest has a steady electric field. Nothing changes, so no
      wave is made.
    \item[\xmark] \textbf{B:} a constant velocity means no acceleration. The electron's fields
      move along with it, but it sends out no wave.
    \item[\cmark] \textbf{C:} the electrons in the antenna move up and down all the time, so
      they are always accelerating. Accelerating charges make changing electric and magnetic
      fields. These keep making each other and travel away as an EM wave, with the same
      frequency as the electrons.
    \item[\xmark] \textbf{D:} a steady current makes a steady magnetic field. The charges move at
      a constant speed, so nothing changes and no wave is made.
  \end{itemize}
\end{practiceans}

\Needspace{12\baselineskip}
\begin{practiceans}
  \correct{B}
  \begin{itemize}
    \item[\xmark] \textbf{A:} $0.20c$ subtracts the two speeds. Speeds never combine with the
      speed of light like this.
    \item[\cmark] \textbf{B:} Einstein's second postulate: every inertial observer measures the
      speed of light in a vacuum as $c$, however fast the source moves. So the station measures
      $1.00c$.
    \item[\xmark] \textbf{C:} $1.67c$ is $c$ divided by $\sqrt{1 - 0.80^2} = 0.60$. That root is
      for time and length, not for the speed of light. Also, nothing travels faster than $c$.
    \item[\xmark] \textbf{D:} $1.80c$ adds the two speeds, as for a ball thrown forward from a
      moving car. This everyday rule fails for light.
  \end{itemize}
\end{practiceans}

\Needspace{13\baselineskip}
\begin{practiceans}
  \correct{A}
  \begin{itemize}
    \item[\cmark] \textbf{A:} Newton's model needs light to go \emph{faster} in water: the water
      pulls the particles in, so the part of their velocity along the normal grows. Foucault
      found the opposite. So Newton's explanation of refraction was wrong. (Huygens' wave model
      predicted slower light, so the result agreed with it.)
    \item[\xmark] \textbf{B:} the idea that light is an EM wave came from Maxwell's theory (1865)
      and Hertz's experiments (1887). A speed in water says nothing about electric and magnetic
      fields.
    \item[\xmark] \textbf{C:} polarisation shows that light is transverse. A speed measurement
      cannot show this.
    \item[\xmark] \textbf{D:} photons come from Einstein (1905) and the photoelectric effect.
      Slower light in water was a prediction of the wave model, not of photons.
  \end{itemize}
\end{practiceans}

\Needspace{15\baselineskip}
\begin{practiceans}
  \correct{D}
  A correct sketch shows three features:
  (1) the hotter curve peaks at a shorter wavelength (Wien's law: \qty{483}{nm} at \qty{6000}{K},
  \qty{725}{nm} at \qty{4000}{K});
  (2) the hotter curve is higher at \emph{every} wavelength, so the curves never cross;
  (3) each curve rises steeply on the short-wavelength side and falls slowly, with a long tail,
  on the long-wavelength side.
  \begin{itemize}
    \item[\xmark] \textbf{A:} the curves are symmetric bells. Real curves are not symmetric
      (feature 3). The curves also cross (feature 2).
    \item[\xmark] \textbf{B:} the \qty{6000}{K} curve peaks at the longer wavelength. The hotter
      body must peak at the shorter wavelength (feature 1).
    \item[\xmark] \textbf{C:} the curves cross. At long wavelengths the \qty{4000}{K} curve is
      above the \qty{6000}{K} curve. A hotter black body gives out more at every wavelength
      (feature 2).
    \item[\cmark] \textbf{D:} shows all three features.
  \end{itemize}
\end{practiceans}

\Needspace{14\baselineskip}
\begin{practiceans}
  \correct{D}
  There are two clues. (1) Every line is at a longer wavelength: a redshift, so P is moving
  \emph{away} from Earth. (2) The lines are wider: P rotates faster, or has a denser atmosphere,
  or both.
  \begin{itemize}
    \item[\xmark] \textbf{A:} moving towards Earth would give a blueshift. A lower density would
      give narrower lines, not wider ones.
    \item[\xmark] \textbf{B:} ``moving away'' is right. But line positions do not tell us the
      temperature. Temperature comes from the peak of the continuous spectrum, or from which
      lines are strong. The question gives neither.
    \item[\xmark] \textbf{C:} moving towards Earth would move the lines to shorter wavelengths
      (a blueshift).
    \item[\cmark] \textbf{D:} matches both clues. Rotation and density both widen lines, so this
      information cannot tell us which one it is.
  \end{itemize}
\end{practiceans}

\Needspace{11\baselineskip}
\begin{practiceans}
  \correct{A}
  The first filter meets unpolarised light, so it passes half:
  $I_1 = \tfrac12 \times \qty{360}{W.m^{-2}} = \qty{180}{W.m^{-2}}$. The second filter meets
  polarised light, so use Malus' law:
  $I_2 = I_1\cos^2\theta = 180 \times (\cos\ang{50.0})^2 = 180 \times 0.4132
  = \qty{74.4}{W.m^{-2}}$.
  \begin{itemize}
    \item[\cmark] \textbf{A:} \qty{74.4}{W.m^{-2}}, as above.
    \item[\xmark] \textbf{B:} \qty{106}{W.m^{-2}} uses $\cos^2\ang{40.0}$. The angle between the
      axes is \ang{50.0}, not $\ang{90} - \ang{50.0}$.
    \item[\xmark] \textbf{C:} \qty{116}{W.m^{-2}} uses $\cos\ang{50.0}$ without squaring:
      $180 \times 0.643$.
    \item[\xmark] \textbf{D:} \qty{149}{W.m^{-2}} forgets the half at the first filter:
      $360 \times 0.413$.
  \end{itemize}
\end{practiceans}

\Needspace{15\baselineskip}
\begin{practiceans}
  \correct{C}
  The same power means that each laser gives out the same energy each second. A red photon has
  less energy than a blue photon, because $E = \dfrac{hc}{\lambda}$ and red light has the longer
  wavelength. So laser R must give out \emph{more} photons each second to carry the same energy:
  $\dfrac{660}{470} = 1.40$ times as many.
  \begin{itemize}
    \item[\xmark] \textbf{A:} wrong way round. A blue photon has more energy: \qty{2.64}{eV},
      compared with \qty{1.88}{eV} for red.
    \item[\xmark] \textbf{B:} all light travels at $c$ in a vacuum (and almost exactly $c$ in
      air), whatever its colour.
    \item[\cmark] \textbf{C:} as explained above.
    \item[\xmark] \textbf{D:} the same power means the same energy each second.
  \end{itemize}
\end{practiceans}

\Needspace{14\baselineskip}
\begin{practiceans}
  \correct{B}
  \textbf{Proper time} is the time between two events, measured by a clock that is present at
  both events. The events are ``leave Earth'' and ``arrive at the station''. Only the crew's
  clock is at both. \textbf{Proper length} is measured in the frame where the two ends are at
  rest. Earth and the station are at rest relative to each other, so observers on Earth measure
  the proper length (the proper distance).
  \begin{itemize}
    \item[\xmark] \textbf{A:} this swaps the two. Earth's clocks are not present at the arrival,
      so they measure a dilated time. For the crew, Earth and the station are moving, so the
      crew measures a contracted distance.
    \item[\cmark] \textbf{B:} as explained above.
    \item[\xmark] \textbf{C:} the crew measures a contracted distance, so it is not the proper
      length.
    \item[\xmark] \textbf{D:} Earth's clocks measure a dilated time, so it is not the proper time.
  \end{itemize}
\end{practiceans}

%% ==================================================================
\Needspace{18\baselineskip}
\subsection*{Section II: short answer}

\Needspace{12\baselineskip}
\begin{practiceans}
  \given{$f = \qty{2.45}{GHz} = \qty{2.45e9}{Hz}$; first to fifth spot: \qty{25.0}{cm}
    $= \qty{0.250}{m}$; neighbouring spots are $\tfrac12\lambda$ apart}
  \find{(a) the speed of light $c$ \quad (b) one improvement, with a reason}
  \formula{spacing $=$ distance $\div$ number of gaps; \ $\lambda = 2 \times \text{spacing}$;
    \ $c = f\lambda$}
  \substitute{Five spots in a line have $5 - 1 = 4$ gaps between them.\par
    \eqlines{\text{spacing} &= \frac{\qty{0.250}{m}}{4} = \qty{0.0625}{m}\\
    \lambda &= 2 \times \qty{0.0625}{m} = \qty{0.125}{m}\\
    c &= \qty{2.45e9}{Hz} \times \qty{0.125}{m} = \qty{3.0625e8}{m.s^{-1}}}}
  \answer{(a) \final{$c = \qty{3.06e8}{m.s^{-1}}$} (3 significant figures).\par
    (b) Any one of these, with its reason:
    \begin{itemize}
      \item Repeat with fresh chocolate several times and average the results. This reduces
            the effect of random errors (better reliability).
      \item Measure across more spots, then divide by the number of gaps. The same reading
            error is then a smaller fraction of the result (better accuracy).
      \item Heat for a shorter time, so the spots are small. Their centres can then be found
            more exactly (better accuracy).
    \end{itemize}}
  \checkstep{the result is $\dfrac{3.06 - 3.00}{3.00}\times 100\,\% = 2.0\,\%$ above the
    accepted value. This is reasonable for spots a few millimetres wide. Using 5 gaps would give
    \qty{2.45e8}{m.s^{-1}}, and forgetting the $\times 2$ would give \qty{1.53e8}{m.s^{-1}}:
    both are far too small.}
  \begin{markingcriteria}
    \criterion{(a) Finds $\lambda = \qty{0.125}{m}$ (4 gaps, then $\times 2$)}{1}
    \criterion{(a) Calculates $c = \qty{3.06e8}{m.s^{-1}}$, with the unit}{1}
    \criterion{(b) Gives a valid change with a reason}{1}
  \end{markingcriteria}
\end{practiceans}

\Needspace{12\baselineskip}
\begin{practiceans}
  \given{$T = 2700 + 273 = \qty{2973}{K}$ (Celsius to kelvin); $b = \qty{2.898e-3}{m.K}$;
    $h = \qty{6.626e-34}{J.s}$; $c = \qty{3.00e8}{m.s^{-1}}$}
  \find{(a) $\lambda_\mathrm{max}$ and the part of the spectrum \quad (b) the photon energy $E$
    in eV}
  \formula{$\lambda_\mathrm{max} = \dfrac{b}{T}$; \quad $E = \dfrac{hc}{\lambda}$, then divide
    by \num{1.602e-19} to change J to eV}
  \substitute{\eqlines{\lambda_\mathrm{max} &= \frac{\qty{2.898e-3}{m.K}}{\qty{2973}{K}}
      = \qty{9.748e-7}{m} = \qty{975}{nm}\\
    E &= \frac{\qty{6.626e-34}{J.s} \times \qty{3.00e8}{m.s^{-1}}}{\qty{9.748e-7}{m}}
      = \qty{2.039e-19}{J}\\
    E &= \frac{\num{2.039e-19}}{\num{1.602e-19}} = \qty{1.273}{eV} \quad\text{(J to eV)}}}
  \answer{(a) \final{$\lambda_\mathrm{max} = \qty{9.75e-7}{m}$} (\qty{975}{nm}). This is in the
    \textbf{infrared}, because it is longer than \qty{700}{nm}. \quad
    (b) \final{$E = \qty{1.27}{eV}$}}
  \checkstep{visible photons have about 1.8 to \qty{3.1}{eV}. An infrared photon should have
    less, and \qty{1.27}{eV} is less \cmark. Using \qty{2700}{K} instead of \qty{2973}{K} gives
    \qty{1073}{nm}: always change to kelvin first.}
  \begin{markingcriteria}
    \criterion{(a) Changes to kelvin and finds $\lambda_\mathrm{max} = \qty{9.75e-7}{m}$}{1}
    \criterion{(a) Names the infrared part of the spectrum}{1}
    \criterion{(b) $E = \qty{1.27}{eV}$}{1}
  \end{markingcriteria}
\end{practiceans}

\Needspace{22\baselineskip}
\begin{practiceans}
  \given{$t_0 = \qty{2.60e-8}{s}$. The pion's own clock is present when it is made and when it
    decays, so this is the proper time. $v = 0.920c = 0.920 \times \qty{3.00e8}{m.s^{-1}}
    = \qty{2.76e8}{m.s^{-1}}$}
  \find{(a) the lifetime $t$ in the laboratory \quad (b) the distance travelled in the
    laboratory}
  \formula{$t = \dfrac{t_0}{\sqrt{\left(1 - \dfrac{v^2}{c^2}\right)}}$; \quad
    distance $=$ speed $\times$ time}
  \substitute{\eqlines{\frac{v^2}{c^2} &= 0.920^2 = 0.8464 \quad\text{(square first)}\\
    \sqrt{1 - 0.8464} &= \sqrt{0.1536} = 0.3919\\
    t &= \frac{\qty{2.60e-8}{s}}{0.3919} = \qty{6.634e-8}{s}\\
    \text{distance} &= \qty{2.76e8}{m.s^{-1}} \times \qty{6.634e-8}{s} = \qty{18.31}{m}}}
  \answer{(a) \final{$t = \qty{6.63e-8}{s}$} \quad (b) \final{\qty{18.3}{m}}}
  \checkstep{$t > t_0$ \cmark. Without time dilation a pion would travel only
    $\qty{2.76e8}{m.s^{-1}} \times \qty{2.60e-8}{s} = \qty{7.18}{m}$. A laboratory distance needs
    the laboratory time.}
  \begin{markingcriteria}
    \criterion{(a) Uses $t_0 = \qty{2.60e-8}{s}$ as the proper time and substitutes correctly}{1}
    \criterion{(a) $t = \qty{6.63e-8}{s}$}{1}
    \criterion{(b) Distance $= \qty{18.3}{m}$, using the laboratory time}{1}
  \end{markingcriteria}
\end{practiceans}

\Needspace{22\baselineskip}
\begin{practiceans}
  \given{unpolarised intensity $I_0$; after P2, $I_2 = 0.30\,I_0$; P1 vertical; P3 horizontal}
  \find{(a) the angle $\theta$ between P1 and P2 \quad (b) $I_3$ as a percentage of $I_0$}
  \formula{first filter: $I_1 = \tfrac12 I_0$; \ each later filter: $I = I_\mathrm{max}\cos^2\theta$,
    with $\theta$ measured from the filter before it}
  \substitute{(a)\quad\eqlines{I_1 &= \tfrac12 I_0 = 0.50\,I_0\\
    0.30\,I_0 &= 0.50\,I_0\cos^2\theta \quad\text{(Malus' law at P2)}\\
    \cos^2\theta &= \frac{0.30}{0.50} = 0.60 \quad\text{(divide both sides by $0.50\,I_0$)}\\
    \cos\theta &= \sqrt{0.60} = 0.7746 \quad\text{(square root)}\\
    \theta &= \cos^{-1}(0.7746) = \ang{39.2}}\par
    (b) P2 is at \ang{39.2} to the vertical. P3 is horizontal, at \ang{90} to the vertical. So the
    angle between P2 and P3 is $\ang{90} - \ang{39.2} = \ang{50.8}$.\par
    \eqlines{I_3 &= 0.30\,I_0 \times (\cos\ang{50.8})^2 = 0.30\,I_0 \times 0.400 = 0.120\,I_0}}
  \answer{(a) \final{$\theta = \ang{39.2}$} \quad (b) \final{12\,\%} of the original intensity}
  \checkstep{P1 and P3 on their own are crossed, so they would pass nothing. P2 between them turns
    the vibration in two steps. The most that can pass is 12.5\,\%, with P2 at \ang{45}. Our
    12\,\% is just below this \cmark.}
  \begin{markingcriteria}
    \criterion{(a) Uses $I_1 = \tfrac12 I_0$ and gets $\cos^2\theta = 0.60$}{1}
    \criterion{(a) $\theta = \ang{39.2}$}{1}
    \criterion{(b) 12\,\% (using the angle \ang{50.8} between P2 and P3)}{1}
  \end{markingcriteria}
\end{practiceans}

\Needspace{22\baselineskip}
\begin{practiceans}
  \given{$d = \qty{0.350}{mm} = \qty{3.50e-4}{m}$; $L = \qty{2.60}{m}$; $m = 1$ to $m = 6$
    bright fringes: \qty{23.5}{mm} $= \qty{2.35e-2}{m}$}
  \find{the wavelength $\lambda$}
  \formula{$\Delta y = \dfrac{\lambda L}{d}$ (small angles). Multiply both sides by $d$, then
    divide both sides by $L$: $\lambda = \dfrac{d\,\Delta y}{L}$}
  \substitute{From $m = 1$ to $m = 6$ there are $6 - 1 = 5$ fringe spacings.\par
    \eqlines{\Delta y &= \frac{\qty{2.35e-2}{m}}{5} = \qty{4.70e-3}{m}\\
    \lambda &= \frac{\qty{3.50e-4}{m} \times \qty{4.70e-3}{m}}{\qty{2.60}{m}}
      = \qty{6.327e-7}{m}}}
  \answer{\final{$\lambda = \qty{6.33e-7}{m} = \qty{633}{nm}$} (red light)}
  \checkstep{the first-order angle is tiny: $\sin\theta = \dfrac{\lambda}{d} = 0.0018$, so
    $\theta = \ang{0.10}$. So the small-angle formula is fine \cmark. Dividing by 6 instead of 5
    would give \qty{527}{nm}.}
  \begin{markingcriteria}
    \criterion{Finds $\Delta y = \qty{4.70}{mm}$ (5 spacings)}{1}
    \criterion{Rearranges to $\lambda = \dfrac{d\,\Delta y}{L}$ and substitutes in metres (or
      an equivalent method, such as $d\sin\theta = m\lambda$ with $y = L\tan\theta$ for both
      fringes)}{1}
    \criterion{$\lambda = \qty{6.33e-7}{m}$ (\qty{633}{nm})}{1}
  \end{markingcriteria}
\end{practiceans}

\Needspace{26\baselineskip}
\begin{practiceans}
  \given{$\lambda = \qty{220}{nm} = \qty{2.20e-7}{m}$; $\phi = \qty{3.66}{eV}$;
    $h = \qty{6.626e-34}{J.s}$; $c = \qty{3.00e8}{m.s^{-1}}$; $q = \qty{1.602e-19}{C}$;
    $m_\mathrm{e} = \qty{9.109e-31}{kg}$}
  \find{(a) the photon energy $E$ \quad (b) $K_\mathrm{max}$ \quad (c) $V_0$ \quad
    (d) $v_\mathrm{max}$}
  \formula{$E = \dfrac{hc}{\lambda}$; \ $K_\mathrm{max} = hf - \phi$; \ $K_\mathrm{max} = qV_0$;
    \ $K = \tfrac12 mv^2$, so $v = \sqrt{\dfrac{2K}{m}}$}
  \substitute{\eqlines{\text{(a) } E &= \frac{\qty{6.626e-34}{J.s} \times \qty{3.00e8}{m.s^{-1}}}
      {\qty{2.20e-7}{m}} = \qty{9.035e-19}{J}\\
    E &= \frac{\num{9.035e-19}}{\num{1.602e-19}} = \qty{5.640}{eV} \quad\text{(J to eV)}\\
    \text{(b) } K_\mathrm{max} &= \qty{5.640}{eV} - \qty{3.66}{eV} = \qty{1.980}{eV}
      \quad\text{(all in eV)}\\
    K_\mathrm{max} &= 1.980 \times \qty{1.602e-19}{J} = \qty{3.172e-19}{J}
      \quad\text{(eV to J)}\\
    \text{(c) } V_0 &= \frac{K_\mathrm{max}}{q} = \frac{\qty{3.172e-19}{J}}{\qty{1.602e-19}{C}}
      = \qty{1.98}{V}\\
    \text{(d) } v_\mathrm{max} &= \sqrt{\frac{2 \times \qty{3.172e-19}{J}}{\qty{9.109e-31}{kg}}}
      = \qty{8.346e5}{m.s^{-1}}}}
  \answer{(a) \final{\qty{5.64}{eV}} \quad (b) \final{\qty{1.98}{eV}} \quad
    (c) \final{$V_0 = \qty{1.98}{V}$}\par
    (d) \final{$v_\mathrm{max} = \qty{8.35e5}{m.s^{-1}}$}}
  \checkstep{\qty{220}{nm} is shorter than magnesium's threshold wavelength,
    $\lambda_0 = \dfrac{hc}{\phi} = \qty{339}{nm}$, so electrons are released \cmark. The speed is
    under 1\,\% of $c$, so $K = \tfrac12 mv^2$ is fine \cmark. The eV number equals the volts
    number \cmark.}
  \begin{markingcriteria}
    \criterion{(a) Photon energy \qty{5.64}{eV}}{1}
    \criterion{(b) $K_\mathrm{max} = \qty{1.98}{eV}$}{1}
    \criterion{(c) $V_0 = \qty{1.98}{V}$}{1}
    \criterion{(d) $v_\mathrm{max} = \qty{8.35e5}{m.s^{-1}}$ (with $K_\mathrm{max}$ in joules)}{1}
  \end{markingcriteria}
\end{practiceans}

\Needspace{12\baselineskip}
\begin{practiceans}
  \given{$N = 420$ lines per mm $= 420 \times 1000 = \num{4.20e5}$ lines per m;
    $\lambda = \qty{656}{nm} = \qty{6.56e-7}{m}$}
  \find{(a) $\theta$ for $m = 1$ \quad (b) the highest order $m$}
  \formula{$d = \dfrac{1}{N}$; \quad $d\sin\theta = m\lambda$, so $\sin\theta = \dfrac{m\lambda}{d}$
    (divide both sides by $d$).\par
    $\sin\theta$ can be at most 1, so $m$ can be at most $\dfrac{d}{\lambda}$.}
  \substitute{\eqlines{d &= \frac{1}{\num{4.20e5}\,\unit{m^{-1}}} = \qty{2.381e-6}{m}\\
    \sin\theta &= \frac{1 \times \qty{6.56e-7}{m}}{\qty{2.381e-6}{m}} = 0.2755\\
    \theta &= \sin^{-1}(0.2755) = \ang{16.0}\\
    \frac{d}{\lambda} &= \frac{\qty{2.381e-6}{m}}{\qty{6.56e-7}{m}} = 3.63}}
  \answer{(a) \final{$\theta = \ang{16.0}$} \quad (b) \final{$m = 3$}: $\dfrac{d}{\lambda} = 3.63$,
    rounded down. For $m = 4$, $\sin\theta$ would have to be
    $\dfrac{4 \times \num{6.56e-7}}{\num{2.381e-6}} = 1.10$, which is impossible.}
  \checkstep{$m = 3$ gives $\sin\theta = 0.827$, so $\theta = \ang{55.7}$, which is less than
    \ang{90} \cmark.}
  \begin{markingcriteria}
    \criterion{(a) $d = \qty{2.38e-6}{m}$ from lines per metre}{1}
    \criterion{(a) $\theta = \ang{16.0}$}{1}
    \criterion{(b) $m = 3$, because $\sin\theta$ cannot be more than 1}{1}
  \end{markingcriteria}
\end{practiceans}

\Needspace{26\baselineskip}
\begin{practiceans}
  \given{each mass $m = \qty{1.884e-28}{kg}$; both are at rest, so all of their mass becomes
    photon energy; $c = \qty{3.00e8}{m.s^{-1}}$; $\qty{1}{MeV} = \qty{1.602e-13}{J}$}
  \find{(a) the total energy $E$ \quad (b) the energy of each photon in MeV \quad (c) why there
    are two photons}
  \formula{$E = mc^2$, with the mass lost $\Delta m = 2m$}
  \substitute{\eqlines{\Delta m &= 2 \times \qty{1.884e-28}{kg} = \qty{3.768e-28}{kg}\\
    E &= \qty{3.768e-28}{kg} \times (\qty{3.00e8}{m.s^{-1}})^2 = \qty{3.391e-11}{J}\\
    E_\mathrm{photon} &= \frac{\qty{3.391e-11}{J}}{2} = \qty{1.696e-11}{J}\\
    E_\mathrm{photon} &= \frac{\num{1.696e-11}}{\num{1.602e-13}}\,\unit{MeV} = \qty{105.9}{MeV}
      \quad\text{(J to MeV)}}}
  \answer{(a) \final{\qty{3.39e-11}{J}} \quad (b) \final{\qty{106}{MeV}} for each photon \quad
    (c) The pair is at rest, so the total momentum before is zero. Momentum is conserved, so the
    total momentum after must also be zero. One photon always carries momentum, so one photon
    cannot do this. Two equal photons moving in opposite directions have a total momentum of
    zero.}
  \checkstep{a muon has about 207 times the mass of an electron. Electron--positron
    annihilation gives two photons of about \qty{0.511}{MeV}, and
    $207 \times \qty{0.511}{MeV} = \qty{106}{MeV}$ \cmark.}
  \begin{markingcriteria}
    \criterion{(a) Total energy \qty{3.39e-11}{J} (uses the mass of both particles)}{1}
    \criterion{(b) \qty{106}{MeV} for each photon}{1}
    \criterion{(c) Explains with conservation of momentum: the total momentum is zero, so two
      photons move in opposite directions}{1}
  \end{markingcriteria}
\end{practiceans}

\Needspace{12\baselineskip}
\begin{practiceans}
  \begin{modelanswer}
  (a) Classical physics (the Rayleigh--Jeans law) predicted that the intensity rises without
  limit as the wavelength gets shorter. So it predicted infinite intensity in the ultraviolet,
  and infinite total energy. This is the ultraviolet catastrophe. Real curves peak and then fall
  to zero.

  (b) A black body contains oscillators: vibrating charges. Planck said that they can only gain
  or lose energy in packets, called quanta. Each quantum has energy $E = hf$. At short wavelengths
  the frequency is high, so each quantum is large. Heat rarely gives one oscillator enough
  energy for even one of these quanta. So very little short-wavelength radiation is given out,
  and the intensity falls to zero. (At long wavelengths the quanta are tiny, so Planck's result
  agrees with classical physics there.)

  (c) Planck quantised only the energy of the oscillators in the walls, and he saw it as a maths
  trick. Einstein said that light \emph{itself} travels as quanta, called photons, each with
  energy $hf$. He used this to explain the photoelectric effect.
  \end{modelanswer}
  \begin{markingcriteria}
    \criterion{(a) States that classical physics predicted the intensity rising without limit
      (to infinity) at short wavelengths}{1}
    \criterion{(b) Links quanta $E = hf$ to large quanta at high frequency, so few are given out
      and the intensity falls to zero}{2}
    \criterion{(b) States Planck's idea that energy is given out in quanta of $hf$, with no link
      to short wavelengths}{1}
    \criterion{(c) States that Einstein quantised light itself (photons), not only the
      oscillators}{1}
  \end{markingcriteria}
\end{practiceans}

\Needspace{12\baselineskip}
\begin{practiceans}
  \given{$m_0 = \qty{9.109e-31}{kg}$ (data sheet); $v = \qty{2.79e8}{m.s^{-1}}$;
    $c = \qty{3.00e8}{m.s^{-1}}$}
  \find{(a) $p_\mathrm{v}$, and how it compares with $m_0v$}
  \formula{$p_\mathrm{v} = \dfrac{m_0 v}{\sqrt{\left(1 - \dfrac{v^2}{c^2}\right)}}$}
  \substitute{\eqlines{\frac{v}{c} &= \frac{\qty{2.79e8}{m.s^{-1}}}{\qty{3.00e8}{m.s^{-1}}}
      = 0.930\\
    \frac{v^2}{c^2} &= 0.930^2 = 0.8649 \quad\text{(square first)}\\
    \sqrt{1 - 0.8649} &= \sqrt{0.1351} = 0.3676\\
    m_0v &= \qty{9.109e-31}{kg} \times \qty{2.79e8}{m.s^{-1}}\\
      &= \qty{2.541e-22}{kg.m.s^{-1}}\\
    p_\mathrm{v} &= \frac{\qty{2.541e-22}{kg.m.s^{-1}}}{0.3676} = \qty{6.914e-22}{kg.m.s^{-1}}}}
  \answer{(a) \final{$p_\mathrm{v} = \qty{6.91e-22}{kg.m.s^{-1}}$}. This is
    $\dfrac{6.914}{2.541} = 2.72$ times $m_0v$. So $p = m_0v$ is far too small at this speed.}
  \checkstep{$p_\mathrm{v}$ is bigger than $m_0v$, as it must be \cmark.}
  \begin{modelanswer}
  (b) As $v$ gets closer to $c$, the root $\sqrt{1 - v^2/c^2}$ gets closer to zero. So the
  momentum $p_\mathrm{v}$, and the energy needed, grow without limit. At $v = c$ they would be
  infinite. A longer accelerator still pushes with a finite force for a finite time. So it can
  only give a finite change in momentum ($\Delta p = F_\mathrm{net}\Delta t$), and finite energy.
  The electron gets closer and closer to $c$, but it can never reach it. Experiments agree:
  electrons given huge energies in accelerators still travel just below $c$.
  \end{modelanswer}
  \begin{markingcriteria}
    \criterion{(a) Correct substitution, with $v$ in \si{m.s^{-1}} and the root}{1}
    \criterion{(a) $p_\mathrm{v} = \qty{6.91e-22}{kg.m.s^{-1}}$, about 2.7 times $m_0v$}{1}
    \criterion{(b) Explains that $p_\mathrm{v}$ (and the energy needed) increases without limit
      as $v$ approaches $c$, and that a finite force, time or energy cannot supply this}{2}
    \criterion{(b) States that the momentum or the energy needed becomes infinite at $c$}{1}
  \end{markingcriteria}
\end{practiceans}

%% ==================================================================
\Needspace{18\baselineskip}
\subsection*{Section III: extended response}

\Needspace{12\baselineskip}
\begin{practiceans}
  \begin{modelanswer}
  \textbf{Before Maxwell.} Some links between electricity and magnetism were known. A current
  makes a magnetic field (Oersted). A changing magnetic field makes an emf (Faraday). But there
  was no single theory, and light was a separate topic.

  \textbf{Unification.} Maxwell wrote all of electricity and magnetism as one set of four laws.
  His own new idea was that a changing \emph{electric} field also makes a magnetic field. So a
  changing electric field makes a changing magnetic field, which makes a changing electric field,
  and so on.

  \textbf{Prediction of EM waves.} This chain keeps itself going. It travels away from
  accelerating charges, even through empty space. So Maxwell predicted electromagnetic waves:
  electric and magnetic fields that oscillate in phase, at \ang{90} to each other and to the
  direction of travel. Nothing limits their frequency, so he predicted a whole spectrum of them.

  \textbf{Prediction of their speed.} His laws give the speed of the waves from two constants
  that are measured with currents and charges:
  \[
    c = \frac{1}{\sqrt{\mu_0\varepsilon_0}}
      = \frac{1}{\sqrt{(4\pi\times10^{-7}) \times (8.854\times10^{-12})}}
      = \frac{1}{\sqrt{1.113\times10^{-17}}} = \qty{3.00e8}{m.s^{-1}}
  \]
  \textbf{What this changed.} This equals the measured speed of light (Fizeau, Foucault). So
  Maxwell concluded that light is an electromagnetic wave. Light was no longer a separate topic:
  visible light is one small band of the EM spectrum, explained by electromagnetism. In 1887
  Hertz made and detected radio waves that behaved as predicted. This confirmed the theory.
  \end{modelanswer}
  \begin{markingcriteria}
    \criterion{Explains how Maxwell unified electricity and magnetism (a changing electric field
      makes a magnetic field); links this to the prediction of self-sustaining EM waves;
      calculates $\dfrac{1}{\sqrt{\mu_0\varepsilon_0}} = \qty{3.00e8}{m.s^{-1}}$ from data-sheet
      values; concludes that light is an EM wave, part of a wider spectrum}{5}
    \criterion{Explains most of these, with a correct calculation, but one link is weak or
      missing}{4}
    \criterion{Describes the unification or the prediction of EM waves, AND gives the
      calculation or the conclusion that light is an EM wave}{3}
    \criterion{Outlines some relevant features (for example, Maxwell predicted EM waves that
      travel at the speed of light)}{2}
    \criterion{Gives one relevant point}{1}
  \end{markingcriteria}
\end{practiceans}

\Needspace{12\baselineskip}
\begin{practiceans}
  \given{$\lambda_\mathrm{A} = \qty{725}{nm} = \qty{7.25e-7}{m}$;
    $\lambda_\mathrm{B} = \qty{322}{nm} = \qty{3.22e-7}{m}$; $b = \qty{2.898e-3}{m.K}$}
  \find{(a) the surface temperatures $T_\mathrm{A}$ and $T_\mathrm{B}$}
  \formula{$\lambda_\mathrm{max} = \dfrac{b}{T}$, rearranged: $T = \dfrac{b}{\lambda_\mathrm{max}}$}
  \substitute{\eqlines{T_\mathrm{A} &= \frac{\qty{2.898e-3}{m.K}}{\qty{7.25e-7}{m}} = \qty{3997}{K}\\
    T_\mathrm{B} &= \frac{\qty{2.898e-3}{m.K}}{\qty{3.22e-7}{m}} = \qty{9000}{K}}}
  \answer{(a) \final{$T_\mathrm{A} = \qty{4.00e3}{K}$} and \final{$T_\mathrm{B} = \qty{9.00e3}{K}$}}
  \checkstep{the Sun (\qty{5800}{K}) peaks near \qty{500}{nm}. A peaks at a longer wavelength, so
    it is cooler (orange-red). B peaks at a shorter wavelength, so it is hotter (blue-white)
    \cmark.}
  \begin{modelanswer}
  (b) \textbf{Composition.} Both stars show all three hydrogen lines, each moved by the same
  small fraction (see below). So both stars contain hydrogen. Star A also shows both sodium
  lines, moved by the same fraction as its hydrogen lines. So A contains sodium. B shows no
  sodium lines. (Which lines are strong also depends on temperature, so this does not prove
  that B has no sodium.)

  \textbf{Motion.} Every line of A is at a shorter wavelength than in the laboratory. Each line
  moves by about 0.03\,\% of its wavelength (for example $656.28 - 656.08 = \qty{0.20}{nm}$).
  This is a blueshift, so A is moving towards Earth. Every line of B is at a longer wavelength,
  by about 0.023\,\% (for example $656.43 - 656.28 = \qty{0.15}{nm}$). This is a redshift, so B
  is moving away from Earth. A's shift is the bigger fraction, so A moves faster along our line of sight.
  (Only motion along the line of sight shows up.)

  \textbf{Rotation or density.} A's lines are narrow, so A rotates slowly and its atmosphere has
  a low density. B's lines are very broad. So B rotates quickly, or its atmosphere is dense, or
  both. The spectrum alone cannot tell us which.
  \end{modelanswer}
  \textit{Extension (not needed for full marks):} the fractional shift equals $\dfrac{v}{c}$. So
  A moves towards Earth at about
  $\num{3.00e8} \times \dfrac{0.20}{656.28} = \qty{9.1e4}{m.s^{-1}}$ (\qty{91}{km.s^{-1}}). This
  formula is not on the formula sheet.
  \begin{markingcriteria}
    \criterion{(a) Both temperatures correct, with units (\qty{4.00e3}{K} and \qty{9.00e3}{K})}{2}
    \criterion{(a) One temperature correct, or both with the same error}{1}
    \criterion{(b) Correctly analyses composition, motion and line width for both stars. Links
      each feature of the data (lines present, direction of shift, line width) to its cause}{4}
    \criterion{(b) Correct conclusions for all three properties, with some links missing}{3}
    \criterion{(b) Correct conclusions for two properties}{2}
    \criterion{(b) One correct conclusion}{1}
  \end{markingcriteria}
\end{practiceans}

\Needspace{12\baselineskip}
\begin{practiceans}
  \begin{modelanswer}
  \textbf{The two models.} Newton (1704) said that light is a stream of tiny particles that
  travel in straight lines. Huygens (1678) said that light is a wave: every point on a wavefront
  sends out small wavelets. He thought the wave was longitudinal, like sound.

  \textbf{Refraction.} Both models explained why light bends towards the normal as it enters
  water, but they needed opposite speeds. In Newton's model the water pulls the particles in.
  The part of their velocity along the normal grows, so light is \emph{faster} in water. In
  Huygens' model the part of the wavefront that enters the water first slows down first. The
  wavefront turns, and light is \emph{slower} in water. In 1850 Foucault found that light is
  slower in water. This ruled out Newton's explanation of refraction and agreed with the wave
  model.

  \textbf{Interference.} In about 1801 Young passed light through two narrow slits and saw bright
  and dark fringes on a screen. Where the path difference is a whole number of wavelengths, the
  waves add (bright). Where it is a whole number plus a half, they cancel (dark). Particles
  cannot cancel each other to make darkness, but waves can. So this was strong evidence for a
  wave model. It also let Young measure the wavelength of light. Fresnel's wave theory and
  Arago's bright spot (1818) added more support.

  \textbf{Polarisation.} A polarising filter passes light that vibrates in one direction only,
  and two crossed filters block all the light. Only transverse waves can be polarised. Newton's
  particles had no natural explanation for this, and Huygens' longitudinal waves cannot be
  polarised. So the wave model had to change: Young and Fresnel proposed that light is a
  \emph{transverse} wave. Maxwell's theory later showed why: light is an electromagnetic wave,
  with fields at \ang{90} to its direction of travel.

  \textbf{Judgement.} Each piece of evidence changed the model. Interference and the slower speed
  in water ruled out the particle model and supported waves. Polarisation showed that the waves
  are transverse. By the late 1800s the evidence strongly supported light as a transverse wave.
  \end{modelanswer}
  \begin{markingcriteria}
    \criterion{Explains how each model accounts for refraction, and links Foucault's result to
      the rejection of Newton's explanation; explains why interference supports a wave model;
      explains why polarisation needs a transverse wave, so Huygens' longitudinal model had to
      change; makes a clear judgement}{6}
    \criterion{Explains all three pieces of evidence and links them to the change of model, but
      one link is weak or missing}{5}
    \criterion{Explains two pieces of evidence and how each changed the model, OR outlines all
      three with weak links}{4}
    \criterion{Describes two pieces of evidence, with some link to the models}{3}
    \criterion{Outlines both models, or describes one piece of evidence}{2}
    \criterion{Gives one relevant point}{1}
  \end{markingcriteria}
\end{practiceans}

% Layout (book QA): 8 lines (not 12), so this answer starts on the page before. Otherwise its
% model answer breaks with only two lines on the next page.
\Needspace{8\baselineskip}
\begin{practiceans}
  \given{$l_0 = \qty{9.00}{km} = \qty{9.00e3}{m}$: the atmosphere is at rest in the Earth frame,
    so this is the proper length. $t_0 = \qty{2.20e-6}{s}$: measured by the muon's own clock, so
    this is the proper time. $v = 0.9960c = 0.9960 \times \qty{3.00e8}{m.s^{-1}}
    = \qty{2.988e8}{m.s^{-1}}$}
  \find{(a) the distance $vt_0$ \quad (b) the lifetime $t$ and the trip time in the Earth frame
    \quad (c) the thickness $l$ in the muon's frame}
  \formula{distance $=$ speed $\times$ time; \
    $t = \dfrac{t_0}{\sqrt{\left(1 - \dfrac{v^2}{c^2}\right)}}$; \
    $l = l_0\sqrt{\left(1 - \dfrac{v^2}{c^2}\right)}$}
  \substitute{First work out the root once:\par
    \eqlines{\sqrt{1 - 0.9960^2} &= \sqrt{1 - 0.992016} = \sqrt{0.007984} = 0.08935}\par
    \eqlines{\text{(a) } vt_0 &= \qty{2.988e8}{m.s^{-1}} \times \qty{2.20e-6}{s} = \qty{657}{m}\\
    \text{(b) } t &= \frac{\qty{2.20e-6}{s}}{0.08935} = \qty{2.462e-5}{s}
      = \qty{24.6}{\micro\second}\\
    \text{trip} &= \frac{\qty{9.00e3}{m}}{\qty{2.988e8}{m.s^{-1}}} = \qty{3.012e-5}{s}
      = \qty{30.1}{\micro\second}\\
    \text{(c) } l &= \qty{9.00e3}{m} \times 0.08935 = \qty{804}{m}}}
  \answer{(a) \final{\qty{657}{m}} \quad (b) \final{$t = \qty{24.6}{\micro\second}$}; trip
    \final{\qty{30.1}{\micro\second}} \quad (c) \final{$l = \qty{804}{m}$}}
  \checkstep{$t > t_0$ and $l < l_0$ \cmark. Now compare the trip time with the lifetime in each
    frame:\par
    \eqlines{\text{Earth frame: }
      \frac{\text{trip}}{\text{lifetime}} &= \frac{\qty{30.1}{\micro\second}}{\qty{24.6}{\micro\second}} = 1.22\\
    \text{muon frame: trip} &= \frac{\qty{804}{m}}{\qty{2.988e8}{m.s^{-1}}}
      = \qty{2.69}{\micro\second}\\
    \frac{\text{trip}}{\text{lifetime}} &= \frac{\qty{2.69}{\micro\second}}{\qty{2.20}{\micro\second}} = 1.22}\par
    The two frames agree \cmark.}
  \begin{modelanswer}
  (d) Without relativity, a muon travels only about \qty{657}{m} in one mean lifetime. The
  \qty{9.00}{km} trip would take $\dfrac{30.1}{2.20} = 13.7$ mean lifetimes, so almost none
  would reach sea level. In the Earth frame, the muon's clock runs slow (time dilation). Its mean
  lifetime is \qty{24.6}{\micro\second}, so the \qty{30.1}{\micro\second} trip lasts only 1.22
  mean lifetimes. In the muon's frame, the atmosphere moves past it, so it is contracted to
  \qty{804}{m} (length contraction). The trip takes only \qty{2.69}{\micro\second}, again 1.22
  mean lifetimes. Both frames predict that a large fraction of the muons survive the trip, and
  this matches the detectors. So the result is evidence for time dilation (Earth frame) and
  length contraction (muon frame).
  \end{modelanswer}
  \begin{markingcriteria}
    \criterion{(a) \qty{657}{m}}{1}
    \criterion{(b) $t = \qty{24.6}{\micro\second}$ and trip time \qty{30.1}{\micro\second}
      (1 mark each)}{2}
    \criterion{(c) $l = \qty{804}{m}$}{1}
    \criterion{(d) Explains that time dilation (Earth frame) and length contraction (muon
      frame) both give a trip of about 1.2 lifetimes, compared with about 14 without relativity,
      so many muons arrive, as observed}{2}
    \criterion{(d) States that time dilation or length contraction lets the muons reach sea
      level}{1}
  \end{markingcriteria}
\end{practiceans}

\Needspace{12\baselineskip}
\begin{practiceans}
  \given{wavelengths in nm; $c = \qty{3.00e8}{m.s^{-1}}$; $q = \qty{1.602e-19}{C}$}
  \find{(a) $f$ for each wavelength \quad (b) the graph \quad (c) $h$ \quad (d) $f_0$, $\phi$ and
    the metal}
  \formula{(a) $c = f\lambda$, so $f = \dfrac{c}{\lambda}$ (divide both sides by $\lambda$).\par
    (c), (d) \eqlines{V_0 &= \frac{h}{q}f - \frac{\phi}{q}}\par
    This is a straight line, with $V_0$ on the vertical axis and $f$ on the horizontal axis.
    So the gradient is $h/q$, which gives $h = q \times \text{gradient}$. The $f$-intercept is
    $f_0$, and $\phi = -q \times (V_0\text{-intercept})$.}
  \substitute{(a) For example, for \qty{240}{nm}:
    $f = \dfrac{\qty{3.00e8}{m.s^{-1}}}{\qty{2.40e-7}{m}} = \qty{1.250e15}{Hz}
    = 12.50 \times 10^{14}$ \si{Hz}. All six:\par
    {\small\begin{tabular}{@{}L{36mm}*{6}{C{10mm}}@{}}
      \toprule
      $\lambda$ (\si{nm}) & 240 & 265 & 290 & 315 & 340 & 365\\
      $f$ ($\times10^{14}$ \si{Hz}) & 12.50 & 11.32 & 10.34 & 9.52 & 8.82 & 8.22\\
      $V_0$ (\si{V}) & 2.33 & 1.77 & 1.43 & 1.05 & 0.81 & 0.52\\
      \bottomrule
    \end{tabular}\par}
    (b) See the graph below. One straight line passes through the middle of the points, with
    points on both sides of it.\par
    (c) Two points far apart \emph{on the line}: (\qty{8.0e14}{Hz}, \qty{0.44}{V}) and
    (\qty{12.0e14}{Hz}, \qty{2.10}{V}).\par
    \eqlines{\text{gradient} &= \frac{\qty{2.10}{V} - \qty{0.44}{V}}{\qty{12.0e14}{Hz} - \qty{8.0e14}{Hz}}\\
      &= \frac{\qty{1.66}{V}}{\qty{4.0e14}{Hz}} = \qty{4.15e-15}{V.s}\\
    h &= \qty{1.602e-19}{C} \times \qty{4.15e-15}{V.s} = \qty{6.648e-34}{J.s}}\par
    (d) The line is $y = mx + c$, so $c = y - mx$. (Here $c$ is the intercept, not the speed of
    light.) Use the point (\qty{8.0e14}{Hz},
    \qty{0.44}{V}):\par
    \eqlines{V_0\text{-intercept} &= \qty{0.44}{V} - \qty{4.15e-15}{V.s} \times \qty{8.0e14}{Hz}\\
      &= \qty{0.44}{V} - \qty{3.32}{V} = \qty{-2.88}{V}}\par
    At the threshold frequency $f_0$, the stopping voltage is zero:\par
    \eqlines{0 &= (\qty{4.15e-15}{V.s})\,f_0 - \qty{2.88}{V}\\
    f_0 &= \frac{\qty{2.88}{V}}{\qty{4.15e-15}{V.s}} = \qty{6.94e14}{Hz}\\
    \phi &= -q \times (\qty{-2.88}{V}) = \qty{4.61e-19}{J} = \qty{2.88}{eV}}}
  \begin{figure}
    \centering
    \begin{tikzpicture}
      \begin{axis}[hsc, width=0.74\linewidth, height=0.36\linewidth,
          xmin=6, xmax=13, ymin=0, ymax=2.6, xtick={6,7,...,13}, ytick={0,0.5,...,2.5},
          yticklabel style={/pgf/number format/.cd, fixed, fixed zerofill, precision=1},
          xlabel={Frequency $f$ ($\times10^{14}$ \si{Hz})},
          ylabel={Stopping voltage $V_0$ (\si{V})}]
        \addplot[hscseriesB, line width=1.1pt, domain=8:12.6, samples=2] {0.415*x - 2.88};
        \addplot[hscseriesB, line width=1pt, densely dotted, domain=6.94:8, samples=2]
          {0.415*x - 2.88};
        \addplot[hscseriesA, only marks, mark=*, mark size=2pt]
          coordinates {(12.50,2.33) (11.32,1.77) (10.34,1.43) (9.52,1.05) (8.82,0.81) (8.22,0.52)};
        \draw[guide] (axis cs:8.0,0.44) -- (axis cs:12.0,0.44) -- (axis cs:12.0,2.10);
        \fill[hscseriesB] (axis cs:8.0,0.44) circle (1.8pt);
        \fill[hscseriesB] (axis cs:12.0,2.10) circle (1.8pt);
        \node[anchor=north, font=\footnotesize] at (axis cs:10.0,0.42) {run $= \qty{4.0e14}{Hz}$};
        \node[anchor=east, font=\footnotesize] at (axis cs:11.92,1.0) {rise $= \qty{1.66}{V}$};
        \fill[hscseriesB] (axis cs:6.94,0) circle (2pt);
        \node[anchor=south west, font=\footnotesize, inner sep=1pt] at (axis cs:6.12,0.72)
          {$f_0 \approx \qty{6.9e14}{Hz}$};
        \draw[->, line width=0.5pt, hscmuted] (axis cs:6.6,0.70) -- (axis cs:6.91,0.08);
      \end{axis}
    \end{tikzpicture}
    \caption*{Model graph: data (dots), line of best fit (dotted past the data), rise and run
      (dashed).}
  \end{figure}
  \answer{(c) \final{$h = \qty{6.65e-34}{J.s}$} \quad (d) \final{$f_0 = \qty{6.94e14}{Hz}$},
    \final{$\phi = \qty{2.88}{eV}$}, so the metal is \textbf{calcium} (\qty{2.87}{eV})}
  \checkstep{$h$ is 0.3\,\% above the data-sheet value \cmark. A hand-drawn line gives slightly
    different values. Accept $h$ from \num{6.4e-34} to \qty{6.9e-34}{J.s}, $f_0$ from \num{6.8e14}
    to \qty{7.1e14}{Hz}, and $\phi$ from 2.8 to \qty{2.95}{eV}. All of these point to calcium.}
  \begin{markingcriteria}
    \criterion{(a) All six frequencies correct}{1}
    \criterion{(b) Axes labelled with quantities and units, a sensible scale, all points plotted
      accurately and one straight line of best fit}{2}
    \criterion{(b) Points plotted with a minor error, or no line of best fit}{1}
    \criterion{(c) Gradient from two points far apart on the line, and
      $h = q \times \text{gradient} \approx \qty{6.6e-34}{J.s}$}{2}
    \criterion{(c) Correct gradient, or correct method with one error}{1}
    \criterion{(d) $f_0 \approx \qty{6.9e14}{Hz}$ and $\phi \approx \qty{2.9}{eV}$, with calcium
      identified}{2}
    \criterion{(d) Correct $f_0$ or correct $\phi$}{1}
  \end{markingcriteria}
\end{practiceans}

\end{practiceanswers}
